% ============================================================================
\clearpage
\section{Digital Signal Processing Fundamentals for Bitwise AI}
\label{sec:appendix-dsp}
% ============================================================================

Throughout this article---particularly in \cref{sec:alt-architectures} and
\cref{sec:scaling-llms}---we noted that state-space models and recurrent
architectures are mathematically identical to classical digital signal
processing (DSP) operations. This appendix provides the DSP background
needed to see those connections clearly, and shows why these operations
are natural candidates for bitwise implementation.

No prior signal processing knowledge is assumed. We build everything from
the simplest possible filter and work toward the exact equations used in
Mamba, RWKV, and Tsetlin Machines.

% ----------------------------------------------------------------------------
\subsection{Why DSP Matters for This Project}
\label{sec:dsp-why}
% ----------------------------------------------------------------------------

The core recurrence of every state-space model is:
%
\begin{equation}
  h_t = \bar{A}\, h_{t-1} + \bar{B}\, x_t
  \label{eq:ssm-recurrence-appendix}
\end{equation}
%
where $h_t$ is the hidden state, $x_t$ is the current input,
$\bar{A}$ is a decay factor, and $\bar{B}$ scales the input.
If you have studied signal processing, you recognise this immediately:
it is a \concept{first-order IIR filter}. If you have not, that term
may mean nothing yet---but it will by the end of this appendix.

The key observation is practical: DSP engineers have implemented
exactly this operation in \emph{fixed-point integer arithmetic} on
hardware as constrained as 8-bit microcontrollers, for applications
ranging from audio equalisers to industrial control loops. They did
not need floating-point units, GPUs, or even integer multipliers in
many cases. The entire field of embedded DSP is, in essence, a
decades-long existence proof that the core operations of modern
neural architectures can run on minimal hardware.

\begin{keyinsight}
The Mamba recurrence is a first-order IIR filter.
The RWKV exponential decay is an exponential moving average.
These are not analogies---they are mathematical identities.
DSP textbooks from the 1970s describe exactly how to implement
them with integer arithmetic and bit shifts.
\end{keyinsight}

% ----------------------------------------------------------------------------
\subsection{IIR Filters (Infinite Impulse Response)}
\label{sec:dsp-iir}
% ----------------------------------------------------------------------------

A \concept{digital filter} takes a sequence of input samples and
produces a sequence of output samples. Filters are classified by
whether their output depends only on current and past inputs, or also
on past outputs.

An \concept{IIR filter} (Infinite Impulse Response) is one whose output
feeds back into itself. The general form is:
%
\begin{equation}
  y[n] = \sum_{k=0}^{M} b_k \, x[n-k]
         \;-\; \sum_{k=1}^{N} a_k \, y[n-k]
  \label{eq:iir-general}
\end{equation}
%
The first sum processes current and past \emph{inputs} (the
\concept{feedforward} path). The second sum uses past \emph{outputs}
(the \concept{feedback} path). It is the feedback that makes the
impulse response ``infinite''---a single input pulse produces output
that, in principle, decays forever.

\subsubsection{First-Order IIR: The Simplest Case}

The simplest possible IIR filter uses one input sample and one past
output:
%
\begin{equation}
  y[n] = b_0 \, x[n] + a_1 \, y[n-1]
  \label{eq:iir-first-order}
\end{equation}
%
Compare this with the SSM recurrence in \cref{eq:ssm-recurrence-appendix}:
they are identical, with $b_0 = \bar{B}$, $a_1 = \bar{A}$,
$x[n] = x_t$, and $y[n] = h_t$.

\begin{figure}[htbp]
\centering
\begin{tikzpicture}[>=Stealth,
  block/.style={draw, rounded corners=3pt, minimum height=0.8cm,
    minimum width=1.4cm, font=\small\sffamily, inner sep=4pt},
  sum/.style={draw, circle, minimum size=0.6cm, font=\sffamily,
    inner sep=0pt},
  arr/.style={-{Stealth[length=4pt]}, thick},
]
  % Input
  \node[font=\sffamily] (xn) at (-2, 0) {$x[n]$};

  % b0 multiply
  \node[block, fill=bitblue!15] (b0) at (0, 0) {$\times\; b_0$};

  % Sum
  \node[sum, fill=bitgreen!15] (plus) at (2.5, 0) {$+$};

  % Output
  \node[font=\sffamily] (yn) at (5, 0) {$y[n]$};

  % Delay
  \node[block, fill=bitorange!15] (delay) at (5, -2) {$z^{-1}$};

  % a1 multiply
  \node[block, fill=bitpurple!15] (a1) at (2.5, -2) {$\times\; a_1$};

  % Connections
  \draw[arr] (xn) -- (b0);
  \draw[arr] (b0) -- (plus);
  \draw[arr] (plus) -- (yn);
  \draw[arr] (yn.south) |- (delay.east);
  \draw[arr] (delay) -- (a1);
  \draw[arr] (a1) -- (plus);

  % Labels
  \node[font=\scriptsize\sffamily, above=2pt, bitblue] at (b0.north)
    {input scaling};
  \node[font=\scriptsize\sffamily, below=2pt, bitorange] at (delay.south)
    {one-sample delay};
  \node[font=\scriptsize\sffamily, below=2pt, bitpurple] at (a1.south)
    {feedback (decay)};
\end{tikzpicture}
\caption{Block diagram of a first-order IIR filter. The delay element
  $z^{-1}$ stores the previous output. This is exactly the structure
  of the Mamba SSM recurrence, with $b_0 = \bar{B}$ and
  $a_1 = \bar{A}$.}
\label{fig:iir-block}
\end{figure}

\subsubsection{Frequency Response and Stability}

Every filter has a \concept{frequency response}---a description of
which frequencies it amplifies and which it attenuates. A first-order
IIR with $0 < a_1 < 1$ acts as a \concept{low-pass filter}: it passes
slow changes (low frequencies) and smooths out rapid changes (high
frequencies). The closer $a_1$ is to~1, the more aggressively it
smooths.

In the SSM context, $\bar{A}$ plays the same role: a large $\bar{A}$
(close to~1) means the state retains a long memory of past inputs,
while a small $\bar{A}$ means the state responds mainly to recent
input. Mamba makes this parameter \emph{input-dependent}, effectively
learning which ``frequencies'' to keep at each time step.

\concept{Stability} requires $|a_1| < 1$. If $|a_1| \geq 1$, the
feedback grows without bound---the filter's output explodes to
infinity. This is the same stability condition that SSMs must satisfy:
the eigenvalues of $\bar{A}$ must lie within the unit circle.

\begin{analogy}
A first-order IIR filter is like a bathtub with a drain that is always
partially open. Water flows in ($x[n]$, the input), and a fraction
leaks out each moment ($a_1 < 1$, the decay). The water level ($y[n]$,
the output) is a smoothed version of the inflow history. If the drain
is fully closed ($a_1 = 1$), the tub overflows---the system is
unstable.
\end{analogy}

% ----------------------------------------------------------------------------
\subsection{Exponential Moving Average}
\label{sec:dsp-ema}
% ----------------------------------------------------------------------------

The \concept{exponential moving average} (EMA) is a specific
first-order IIR filter that appears throughout machine learning,
often without being identified as a filter. Its equation is:
%
\begin{equation}
  y[n] = \alpha \, x[n] + (1 - \alpha) \, y[n-1],
  \quad 0 < \alpha < 1
  \label{eq:ema}
\end{equation}
%
This is \cref{eq:iir-first-order} with $b_0 = \alpha$ and
$a_1 = (1 - \alpha)$. Because $b_0 + a_1 = 1$, the EMA preserves
the DC gain (a constant input produces the same constant output),
making it a natural choice for computing running averages.

\subsubsection{How the Weights Decay}

Expanding the recursion reveals the weighting pattern:
%
\begin{align}
  y[n] &= \alpha \, x[n]
         + \alpha(1-\alpha) \, x[n-1]
         + \alpha(1-\alpha)^2 \, x[n-2] + \cdots
  \label{eq:ema-expanded}
\end{align}
%
Each past sample is weighted by $\alpha(1-\alpha)^k$, which decays
exponentially with age~$k$. The parameter $\alpha$ controls how
quickly old data is forgotten: a large $\alpha$ (close to~1) tracks
the input closely with little smoothing; a small $\alpha$ produces
heavy smoothing with long memory.

\begin{figure}[htbp]
\centering
\begin{tikzpicture}
\begin{axis}[
  width=10cm, height=5cm,
  xlabel={Time step $k$ (samples ago)},
  ylabel={Weight $\alpha(1-\alpha)^k$},
  ymin=0, ymax=0.55,
  xmin=0, xmax=15,
  xtick={0,2,...,14},
  legend style={at={(0.97,0.97)}, anchor=north east,
    font=\small, draw=none, fill=white, fill opacity=0.8},
  grid=major,
  grid style={gray!20},
]
  % alpha = 0.5
  \addplot[bitblue, thick, mark=*, mark size=2pt]
    coordinates {
      (0,0.5) (1,0.25) (2,0.125) (3,0.0625)
      (4,0.03125) (5,0.015625) (6,0.0078125) (7,0.00390625)
      (8,0.00195) (9,0.000977) (10,0.000488) (11,0.000244)
      (12,0.000122) (13,0.0000610) (14,0.0000305)
    };
  \addlegendentry{$\alpha = 0.5$ (fast decay)}

  % alpha = 0.2
  \addplot[bitorange, thick, mark=square*, mark size=2pt]
    coordinates {
      (0,0.2) (1,0.16) (2,0.128) (3,0.1024)
      (4,0.08192) (5,0.065536) (6,0.0524) (7,0.04194)
      (8,0.03355) (9,0.02684) (10,0.02147) (11,0.01718)
      (12,0.01374) (13,0.01100) (14,0.00879)
    };
  \addlegendentry{$\alpha = 0.2$ (slow decay)}

  % alpha = 0.1
  \addplot[bitgreen, thick, mark=triangle*, mark size=2.5pt]
    coordinates {
      (0,0.1) (1,0.09) (2,0.081) (3,0.0729)
      (4,0.06561) (5,0.05905) (6,0.05314) (7,0.04783)
      (8,0.04305) (9,0.03874) (10,0.03487) (11,0.03138)
      (12,0.02824) (13,0.02542) (14,0.02288)
    };
  \addlegendentry{$\alpha = 0.1$ (very slow decay)}
\end{axis}
\end{tikzpicture}
\caption{Exponential decay of sample weights in the EMA for different
  values of $\alpha$. A large $\alpha$ concentrates weight on the most
  recent sample; a small $\alpha$ spreads weight over a longer history.
  The RWKV architecture learns this decay rate per channel.}
\label{fig:ema-decay}
\end{figure}

\begin{figure}[htbp]
\centering
\begin{tikzpicture}
\begin{axis}[
  width=10cm, height=5cm,
  xlabel={Time step $n$},
  ylabel={Value},
  xmin=0, xmax=30,
  legend style={at={(0.97,0.97)}, anchor=north east,
    font=\small, draw=none, fill=white, fill opacity=0.8},
  grid=major,
  grid style={gray!20},
]
  % Raw signal (noisy sine)
  \addplot[bitgray, thin, mark=none]
    coordinates {
      (0,2.1) (1,3.8) (2,1.2) (3,4.5) (4,2.9) (5,5.1) (6,3.3)
      (7,5.8) (8,4.1) (9,6.2) (10,4.8) (11,5.5) (12,4.2) (13,5.0)
      (14,3.8) (15,4.1) (16,3.0) (17,3.5) (18,2.2) (19,2.8)
      (20,1.5) (21,2.1) (22,0.8) (23,1.5) (24,0.5) (25,1.2)
      (26,0.3) (27,1.8) (28,0.9) (29,2.5) (30,1.1)
    };
  \addlegendentry{Raw signal $x[n]$}

  % EMA output (alpha=0.3)
  \addplot[bitblue, very thick, mark=none]
    coordinates {
      (0,2.1) (1,2.61) (2,2.19) (3,2.88) (4,2.89) (5,3.55)
      (6,3.48) (7,4.17) (8,4.15) (9,4.77) (10,4.78) (11,5.00)
      (12,4.76) (13,4.83) (14,4.52) (15,4.40) (16,3.98)
      (17,3.84) (18,3.35) (19,3.18) (20,2.68) (21,2.50)
      (22,1.99) (23,1.85) (24,1.44) (25,1.37) (26,1.05)
      (27,1.27) (28,1.16) (29,1.56) (30,1.42)
    };
  \addlegendentry{EMA output ($\alpha = 0.3$)}
\end{axis}
\end{tikzpicture}
\caption{An EMA smooths a noisy signal by weighting recent values
  more heavily than older ones. This is the same operation performed
  by batch normalisation's running statistics and by RWKV's
  time-mixing mechanism.}
\label{fig:ema-smoothing}
\end{figure}

\subsubsection{Where EMAs Appear in Neural Networks}

The EMA is ubiquitous in deep learning:

\begin{itemize}[leftmargin=2em]
  \item \textbf{Batch normalisation} maintains running mean and
    variance using an EMA with a typical momentum of $\alpha = 0.1$.
  \item \textbf{Optimiser momentum} (SGD with momentum, Adam's first
    moment) is an EMA of the gradient.
  \item \textbf{RWKV's time-mixing} uses a learned exponential decay
    $e^{-w}$ as the forgetting factor, which is an EMA with
    $\alpha = 1 - e^{-w}$.
  \item \textbf{RetNet's multi-scale retention} uses different decay
    rates per head---multiple EMAs running in parallel, each tuned
    to a different time scale.
\end{itemize}

\begin{engineernote}
If you have ever implemented a simple moving average in firmware or a
running-statistics tracker in Python, you have already used the exact
same recurrence that powers modern sequence models. The only
difference is that models like Mamba learn the filter coefficients
from data, while traditional DSP designs them by hand.
\end{engineernote}

% ----------------------------------------------------------------------------
\subsection{Leaky Integrators}
\label{sec:dsp-leaky}
% ----------------------------------------------------------------------------

A \concept{leaky integrator} is yet another name for the same
first-order IIR structure, written with a different convention:
%
\begin{equation}
  y[n] = (1 - \lambda) \, y[n-1] + \lambda \, x[n],
  \quad 0 < \lambda < 1
  \label{eq:leaky-integrator}
\end{equation}
%
This is identical to the EMA in \cref{eq:ema} with $\lambda = \alpha$.
The name ``leaky integrator'' comes from neuroscience and control
theory, where it models systems that accumulate input but gradually
lose their stored value.

\subsubsection{The Leaky Bucket}

\begin{analogy}
Imagine a bucket with a small hole in the bottom. You pour water in
at a rate proportional to the input signal $x[n]$. Between pours,
water leaks out through the hole at a rate proportional to the current
water level. The water level $y[n]$ is a smoothed, slightly delayed
version of the inflow---it rises when input is sustained and falls
when input stops. The hole size controls the ``memory'': a large hole
(large $\lambda$) tracks input quickly but forgets fast; a small hole
(small $\lambda$) responds slowly but remembers longer.
\end{analogy}

\begin{figure}[htbp]
\centering
\begin{tikzpicture}[>=Stealth]
  % Bucket
  \draw[thick, bitblue] (0,0) -- (0,3) -- (3,3) -- (3,0);

  % Water level
  \fill[bitblue!20] (0.05,0.05) rectangle (2.95,1.8);
  \node[font=\sffamily, bitblue] at (1.5,1.0) {$y[n]$};

  % Input arrow
  \draw[-{Stealth[length=5pt]}, very thick, bitgreen]
    (1.5,4.2) -- (1.5,3.2);
  \node[font=\sffamily\bfseries, bitgreen, above] at (1.5,4.2)
    {$\lambda \cdot x[n]$};
  \node[font=\scriptsize\sffamily, bitgreen] at (1.5,3.7) {input};

  % Leak arrow
  \draw[-{Stealth[length=5pt]}, very thick, bitorange]
    (3.2,0.5) -- (4.5,0.5);
  \node[font=\sffamily\bfseries, bitorange, right] at (4.5,0.5)
    {$\lambda \cdot y[n{-}1]$};
  \node[font=\scriptsize\sffamily, bitorange] at (5.5,0.1) {leak};

  % Hole
  \fill[white] (2.85,0.2) rectangle (3.15,0.8);
  \draw[thick, bitorange] (3,0.2) -- (3,0.8);

  % Bottom
  \draw[thick, bitblue] (0,0) -- (3,0);

  % Labels
  \node[font=\scriptsize\sffamily, bitgray, align=center] at (1.5,-0.7)
    {Small hole $\to$ slow leak $\to$ long memory\\
     Large hole $\to$ fast leak $\to$ short memory};
\end{tikzpicture}
\caption{The leaky bucket analogy for a leaky integrator. Water (signal)
  pours in and gradually leaks out. The water level represents the
  accumulated state. Tsetlin Machine automata use a similar
  accumulate-and-threshold mechanism---see \cref{sec:tsetlin}.}
\label{fig:leaky-bucket}
\end{figure}

\subsubsection{Connection to Tsetlin Machines}

The Tsetlin Machine automaton (described in \cref{sec:tsetlin})
is structurally similar to a leaky integrator with a threshold.
Each automaton has a state counter that increments or decrements
based on feedback---this is the accumulation step. The counter
is bounded (it cannot go below 1 or above $2N$), which acts as
a clamp rather than a continuous leak. But the principle is the
same: accumulate evidence over time, and make a binary decision
(include or exclude a literal) based on whether the accumulated
state crosses a threshold.

\begin{keyinsight}
The leaky integrator, the EMA, and the first-order IIR filter are
all the same equation with different names, used by different
communities. Recognising this unifies DSP, neuroscience, control
theory, and modern sequence modelling under one framework.
\end{keyinsight}

% ----------------------------------------------------------------------------
\subsection{Bitwise Implementation of DSP Operations}
\label{sec:dsp-bitwise}
% ----------------------------------------------------------------------------

This is the section that ties DSP directly to the goals of this
project. We show that every operation discussed above can be
implemented without floating-point arithmetic---and in many cases,
without any multiplication at all.

\subsubsection{Fixed-Point Arithmetic}

\concept{Fixed-point arithmetic} represents fractional numbers as
integers with an implied denominator. Instead of storing $0.75$ as a
float, we store it as $192$ with the understanding that the true value
is $192 / 256$. The denominator is always a power of two
(here~$2^8 = 256$), so we say the number uses \emph{Q8 format}---8
fractional bits.

\begin{engineernote}
In Q8 format, the value $0.75$ is stored as the integer
$\lfloor 0.75 \times 256 \rfloor = 192 = \texttt{0xC0}$.
To multiply $x$ by this fixed-point $0.75$:
compute $x \times 192$, then shift right by~8.
The shift replaces the division by $256$.
The result is an integer approximation of $0.75 \times x$.
\end{engineernote}

Multiplication of two Q8 numbers produces a Q16 result (the fractional
bits add), so a right shift of~8 after multiplication restores the
original format. This is one integer multiply and one shift---no
floating-point unit required.

\subsubsection{Power-of-Two Coefficients: Eliminating Multiplication}

If we constrain the filter coefficient to be a power of two, the
multiplication disappears entirely. If $\alpha = 1/2^k$, then:
%
\begin{equation}
  \alpha \cdot x = \frac{x}{2^k} = x \gg k
  \label{eq:shift-multiply}
\end{equation}
%
where $\gg$ denotes a right bit shift. A bit shift executes in a
single clock cycle and consumes negligible energy compared to even
an integer multiplication.

Common useful values:
\begin{itemize}[leftmargin=2em]
  \item $\alpha = 1/2 \;(k=1)$: $\alpha \cdot x = x \gg 1$
  \item $\alpha = 1/4 \;(k=2)$: $\alpha \cdot x = x \gg 2$
  \item $\alpha = 1/8 \;(k=3)$: $\alpha \cdot x = x \gg 3$
  \item $\alpha = 3/4$: $(x \gg 1) + (x \gg 2)$ (two shifts + one add)
\end{itemize}

\subsubsection{The EMA with Bit Shifts Only}

The EMA equation $y[n] = \alpha \, x[n] + (1-\alpha) \, y[n-1]$
can be rearranged as:
%
\begin{equation}
  y[n] = y[n-1] + \alpha \bigl(x[n] - y[n-1]\bigr)
  \label{eq:ema-rearranged}
\end{equation}
%
With $\alpha = 1/2^k$, this becomes:
%
\begin{equation}
  y[n] = y[n{-}1] + \bigl(x[n] - y[n{-}1]\bigr) \gg k
  \label{eq:ema-bitshift}
\end{equation}
%
This is \emph{one subtraction, one shift, and one addition}. No
multiplication whatsoever. This is how embedded systems engineers
have implemented running averages on 8-bit microcontrollers since
the 1980s.

\begin{figure}[htbp]
\centering
\begin{tikzpicture}[>=Stealth,
  block/.style={draw, rounded corners=3pt, minimum height=0.7cm,
    minimum width=1.6cm, font=\small\sffamily, inner sep=3pt},
  sum/.style={draw, circle, minimum size=0.5cm, font=\sffamily,
    inner sep=0pt},
  arr/.style={-{Stealth[length=4pt]}, thick},
]
  % Input
  \node[font=\sffamily] (xn) at (-2.5, 0) {$x[n]$};

  % Subtract
  \node[sum, fill=bitorange!15] (sub) at (-0.3, 0) {$-$};

  % Shift
  \node[block, fill=bitblue!15] (shift) at (2, 0) {$\gg k$};

  % Add
  \node[sum, fill=bitgreen!15] (add) at (4.5, 0) {$+$};

  % Output
  \node[font=\sffamily] (yn) at (7, 0) {$y[n]$};

  % Delay
  \node[block, fill=bitpurple!15] (delay) at (7, -1.8) {$z^{-1}$};

  % Connections
  \draw[arr] (xn) -- (sub);
  \draw[arr] (sub) -- node[above, font=\scriptsize\sffamily] {$x[n]-y[n{-}1]$} (shift);
  \draw[arr] (shift) -- (add);
  \draw[arr] (add) -- (yn);
  \draw[arr] (yn.south) |- (delay.east);
  \draw[arr] (delay.west) -| (sub.south);
  \draw[arr] (delay.west) -| (add.south);

  % Annotation
  \node[font=\scriptsize\sffamily, bitblue, above=2pt] at (shift.north)
    {bit shift (no multiply!)};
\end{tikzpicture}
\caption{Bit-shift-only EMA implementation. The entire filter uses one
  subtraction, one bit shift, and one addition---no multiplier needed.
  This structure implements the same recurrence used by RWKV and Mamba.}
\label{fig:ema-bitshift}
\end{figure}

\subsubsection{Cascaded First-Order Sections}

Complex filters with sharp frequency responses are traditionally
built from higher-order equations with many coefficients. An
alternative is to \concept{cascade} multiple first-order sections in
series: the output of one filter feeds into the input of the next.
Each section is a simple bit-shift EMA, but the cascade achieves
steeper roll-off.

This is directly relevant to deep neural networks: stacking multiple
SSM layers is equivalent to cascading multiple IIR filter stages.
Each layer is a simple first-order recurrence, but the depth of the
network gives it the expressive power of a high-order filter.

\subsubsection{Quantisation Effects}

Fixed-point and bit-shift implementations introduce
\concept{quantisation noise}---rounding errors from discarding bits
during shifts. For a Q$n$ representation, the worst-case error per
operation is $2^{-n}/2$. Over many cascaded operations, these errors
accumulate.

Strategies for managing quantisation noise include:

\begin{itemize}[leftmargin=2em]
  \item \textbf{Guard bits}: use a wider internal representation
    (e.g., Q16 internally, Q8 at input/output) to keep precision
    during intermediate computations.
  \item \textbf{Rounding instead of truncation}: adding
    $2^{k-1}$ before shifting right by $k$ rounds to the nearest
    integer instead of always rounding down.
  \item \textbf{Dithering}: adding a small random value before
    quantisation to break up systematic rounding patterns.
  \item \textbf{Noise shaping}: designing the filter to push
    quantisation noise into frequency bands where it matters less.
\end{itemize}

The neural network analogue of quantisation noise management is the
\concept{straight-through estimator} (STE) discussed in
\cref{sec:bnn}: both address the fundamental challenge of maintaining
gradient flow (or signal quality) through discretisation steps.

\subsubsection{Worked Example: Integer-Only EMA}

The following pseudocode implements a complete EMA filter using only
integer arithmetic and bit shifts, suitable for an 8-bit
microcontroller or a bitwise neural network inference engine.

\begin{algorithm}[htbp]
\caption{Integer-only EMA with bit shifts}
\label{alg:ema-integer}
\begin{algorithmic}[1]
\State \textbf{Parameters:} shift $k$ (sets $\alpha = 1/2^k$)
\State \textbf{State:} accumulator $y \gets 0$ \Comment{Q8 fixed-point}
\Procedure{EmaUpdate}{$x$} \Comment{$x$ is Q8 input}
  \State $\text{diff} \gets x - y$ \Comment{subtraction}
  \State $\text{step} \gets \text{diff} \gg k$ \Comment{bit shift (= $\alpha \cdot \text{diff}$)}
  \State $y \gets y + \text{step}$ \Comment{addition}
  \State \Return $y$
\EndProcedure
\end{algorithmic}
\end{algorithm}

Three operations. No floating point. No multiplication. This is the
core recurrence of a state-space model, implemented in the simplest
possible arithmetic.

\subsubsection{Cost Comparison}

\Cref{tab:dsp-cost} summarises the computational cost of the key
operations across three implementation strategies. The energy figures
are approximate and based on published estimates for 45\,nm CMOS
technology.

\begin{table}[htbp]
\centering
\caption{Cost comparison for DSP/SSM operations across implementation
  strategies. Energy estimates are for 45\,nm CMOS at representative
  bit widths (FP32, INT16, INT8).}
\label{tab:dsp-cost}
\begin{tabularx}{\textwidth}{@{}l >{\centering\arraybackslash}X
  >{\centering\arraybackslash}X >{\centering\arraybackslash}X@{}}
  \toprule
  \textbf{Operation}
    & \textbf{Floating-point}
    & \textbf{Fixed-point (INT)}
    & \textbf{Bit-shift only} \\
  \midrule
  Multiply by $\alpha$
    & 1 FP multiply
    & 1 INT multiply
    & 1 shift \\
  EMA update
    & 2 FP mul + 1 add
    & 2 INT mul + 1 add
    & 1 sub + 1 shift + 1 add \\
  Approx.\ energy per op
    & ${\sim}100$\,pJ
    & ${\sim}10$\,pJ
    & ${\sim}1$\,pJ \\
  Hardware required
    & FPU
    & Integer ALU
    & Barrel shifter + adder \\
  Precision
    & ${\sim}7$ decimal digits
    & Configurable (Q$n$)
    & Power-of-two $\alpha$ only \\
  \bottomrule
\end{tabularx}
\end{table}

The two orders-of-magnitude reduction from floating-point to bit-shift
operations is the fundamental motivation for this entire project. If
the core operations of modern sequence models \emph{are} DSP
operations, and DSP operations \emph{can} be implemented with bit
shifts, then the path from transformer-class architectures to
bitwise-only hardware is not a speculative leap---it is a well-trodden
engineering path.

% ----------------------------------------------------------------------------
\subsection{Connection to Neural Network Building Blocks}
\label{sec:dsp-nn-connection}
% ----------------------------------------------------------------------------

\Cref{tab:dsp-nn-mapping} maps DSP concepts to their neural network
counterparts. These mappings are not analogies---they are mathematical
identities or direct structural correspondences.

\begin{table}[htbp]
\centering
\caption{DSP--neural network correspondence. Each row shows a DSP
  operation and the neural network component that implements the same
  (or structurally equivalent) mathematics.}
\label{tab:dsp-nn-mapping}
\begin{tabularx}{\textwidth}{@{}l >{\raggedright\arraybackslash}X
  >{\raggedright\arraybackslash}X@{}}
  \toprule
  \textbf{DSP concept}
    & \textbf{Neural network equivalent}
    & \textbf{Relationship} \\
  \midrule
  First-order IIR filter
    & SSM recurrence (Mamba)
    & Mathematical identity \\
  Exponential moving average
    & Exponential decay attention (RWKV)
    & Mathematical identity \\
  Leaky integrator
    & Tsetlin automaton state update
    & Structural analogue \\
  FIR filter (finite impulse response)
    & Conv1d layer
    & Mathematical identity \\
  Delay line ($z^{-1}$)
    & Causal masking / positional encoding
    & Functional equivalence \\
  Cascaded filter sections
    & Stacked SSM layers (deep network)
    & Structural analogue \\
  Quantisation noise
    & Weight/activation quantisation error
    & Same phenomenon \\
  Noise shaping
    & Straight-through estimator
    & Analogous strategy \\
  \bottomrule
\end{tabularx}
\end{table}

\begin{keyinsight}
These are not analogies---they are mathematical identities. The SSM
recurrence $h_t = \bar{A} h_{t-1} + \bar{B} x_t$ \textbf{is} a
first-order IIR filter. The RWKV exponential decay \textbf{is} an
exponential moving average. A Conv1d layer \textbf{is} an FIR filter.
Recognising this means that sixty years of DSP implementation
expertise---fixed-point arithmetic, bit-shift filters, cascaded
sections, noise shaping---transfers directly to neural network
hardware design.
\end{keyinsight}

This connection closes the loop between the four bitwise paradigms
surveyed in this article and the sequence modelling architectures
discussed in \cref{sec:alt-architectures}. The operations at the heart
of modern sequence models are the same operations that DSP engineers
have been implementing in integer and bitwise arithmetic for decades.
The remaining challenge is not whether these operations \emph{can} be
made bitwise---it is whether the full \emph{system} (training, learned
coefficients, and all) can achieve competitive accuracy when confined
to bitwise-only inference.
